\clearpage
\VerdierPhysicalPageAnchor{0076}%
\VerdierSemanticAnchor{ch1-2-4-2-countable-products-sums}
现在假设 $C$ 中的可数乘积（相应地，可数直和）可表示。规定：

\VerdierSemanticAnchor{ch1-2-4-2-eq6}
\begin{equation}
  \begin{aligned}
    \operatorname{comp}_{\Pi}F
      &=\int_n^{\Pi}F^{[n]},\\
    \text{（相应地，}\operatorname{comp}_{\Sigma}F
      &=\int_n^{\Sigma}F^{[n]}\text{）}.
  \end{aligned}
  \tag{2.4.2.6}
\end{equation}

\VerdierSemanticAnchor{ch1-2-4-2-pi-sigma-extension}
函子
$\operatorname{comp}_{\Pi}F$（相应地，$\operatorname{comp}_{\Sigma}F$）：
$\prod_{i\in[n]}\operatorname{comp}(C_i)\longrightarrow
\operatorname{comp}(C)$ 称为函子 $F$ 在复形上的 $\Pi$-扩张
（相应地，$\Sigma$-扩张）。

\VerdierSemanticAnchor{ch1-2-4-2-koszul-twist}
由于 $F$ 是严格 $\mathbb Z[n]$-函子，而函子
$\int_n^{\Pi}$（相应地，$\int_n^{\Sigma}$）是由
$\mathbb Z[n]$ 的有序基 $e_1,\ldots,e_n$ 所对应的 Koszul 余循环扭曲的
$\mathbb Z[n]$-函子，函子
$\operatorname{comp}_{\Pi}F$（相应地，$\operatorname{comp}_{\Sigma}F$）
也是由同一有序基所对应的 Koszul 余循环扭曲的
$\mathbb Z[n]$-函子，这与符号规则 I (1.6.6) 一致。因此，对每个
$i\in[n]$，都有典范同构：

\VerdierSemanticAnchor{ch1-2-4-2-eq7}
\begin{equation}
  \begin{aligned}
    \operatorname{comp}_{\Pi}F(X_1,\ldots,X_i[1],\ldots,X_n)
      &\xrightarrow{\sim}
      \operatorname{comp}_{\Pi}F(X_1,\ldots,X_i,\ldots,X_n)[1],\\
    \text{（相应地，}\operatorname{comp}_{\Sigma}F(X_1,\ldots,X_i[1],\ldots,X_n)
      &\xrightarrow{\sim}
      \operatorname{comp}_{\Sigma}F(X_1,\ldots,X_i,\ldots,X_n)[1]\text{）}
  \end{aligned}
  \tag{2.4.2.7}
\end{equation}

\VerdierSemanticAnchor{ch1-2-4-2-canonical-pi}
并将它们记为 $p_i$。

\VerdierSemanticAnchor{ch1-2-4-2-anticommutativity}
这些同构具有 (1.6.6.3) 所述的反交换性质。

\VerdierSemanticAnchor{ch1-2-4-2-summable-restriction-equality}
函子 $\operatorname{comp}_{\Pi}F$ 与
$\operatorname{comp}_{\Sigma}F$ 在
$\prod_{i\in[n]}\operatorname{comp}(C_i)$ 的可求和对象上的限制相等；
这一限制总能定义，而无须对范畴 $C$ 作任何假设。

\VerdierSemanticAnchor{ch1-2-4-3}
\subsection*{2.4.3.}
现设
$F,F'\colon\prod_{i\in[n]}C_i\longrightarrow C$ 为两个多加法函子，
$h\colon F\longrightarrow F'$ 为一个函子态射。由此得到一个函子态射：

\VerdierSemanticAnchor{ch1-2-4-3-induced-morphism}
\begin{equation*}
  h^{[n]}\colon F^{[n]}\longrightarrow F'^{[n]},
\end{equation*}
其定义为：

\VerdierSemanticAnchor{ch1-2-4-3-eq1}
\begin{equation}
  \begin{aligned}
    h^{[n]}(X_1,\ldots,X_n)^{\sigma}
      &=h(X_1^{l_1},\ldots,X_n^{l_n}),\\
    \sigma&=\sum_{i\in[n]}l_i e_i\in\mathbb Z[n],\\
    (X_1,\ldots,X_n)&\in
      \operatorname{Ob}\!\left(\prod_{i\in[n]}\operatorname{comp}(C_i)\right).
  \end{aligned}
  \tag{2.4.3.1}
\end{equation}
